\documentclass{article}
\usepackage[utf8]{inputenc}
\usepackage[T1]{fontenc}
\usepackage[portuguese]{babel}
\usepackage{amsmath}
\usepackage{amsfonts}
\usepackage{amssymb}
\usepackage{graphicx}
\usepackage{geometry} % Para ajustar as margens
\geometry{a4paper, top=2.5cm, bottom=2.5cm, left=3cm, right=2cm} % Margens ABNT-ish

% Remover cabeçalhos/rodapés e numeração de página da capa
\pagestyle{empty}

% --- Informações da Capa ---
% O título principal da capa
\title{\textbf{TRABALHO: ANÁLISE DE RETINOPATIA DIABÉTICA}}

% O autor
\author{João Henrique Serodio Ulbinski}

% A data do trabalho
\date{2025}

\begin{document}

% --- Conteúdo da Capa ---

% Centraliza todo o conteúdo da capa
\begin{center}
    % Espaço vertical para posicionamento
    \vspace*{\fill} % Empurra o conteúdo para o meio/abaixo da página

    % Nome da instituição
    \Large{\textbf{UNICENTRO - Campus CEDETEG Guarapuava}} \\
    \vspace{0.5cm}

    % Curso e Ano
    \normalsize
    Curso: Matemática Aplicada Computacional \\
    Segundo Ano \\
    \vspace{1cm}

    % Título principal do trabalho (definido com \title)
    \maketitle

    % Disciplina
    \vspace{1cm}
    Disciplina: Probabilidade \\
    \vspace{1cm}

    % Nome do Aluno (definido com \author)
    \normalsize
    Aluno: João Henrique Serodio Ulbinski \\
    \vspace*{\fill} % Empurra o conteúdo para o final da página

    % Ano (definido com \date, mas reescrito para maior controle)
    \normalsize
    Guarapuava \\
    2025

\end{center}

\end{document}